\documentclass{beamer}

\usetheme{Berlin}
\usecolortheme{Orchid}
\useoutertheme{miniframes}
\setbeamertemplate{navigation symbols}{}

\usepackage[utf8]{inputenc}
\usepackage[T1]{fontenc}
\usepackage{amsmath}
\usepackage{amssymb}
\usepackage{booktabs}
\usepackage{graphicx}
\graphicspath{{../images/}}

\usepackage{xcolor}
\usepackage{listings}
\usepackage{hyperref}
\usepackage{tikz}
\usetikzlibrary{arrows.meta,positioning,shapes.geometric,calc}

\lstset{
  language=Java,
  basicstyle=\ttfamily\small,
  keywordstyle=\color{blue},
  commentstyle=\color{gray},
  stringstyle=\color{teal},
  showstringspaces=false,
  columns=fullflexible,
  breaklines=true
}

% Per-slide source line (references shown on the slide itself).
\newcommand{\slidesources}[1]{%
  \vfill
  {\tiny\color{gray}\raggedright\sloppy\parbox{\linewidth}{Sources: #1}\par}%
}

% Unit-wise source strings for this subject (used by slides/notes).
% Path is relative to the lecture `latex/` folder (the compilation CWD).
% OOPS with Java (CSEG1044) — unit-wise sources
%
% These are short, student-facing reference strings intended to be used with:
%   - \slidesources{...}  (in beamer slides)
%   - \notesources{...}   (in lecture notes)
%
% Style inspiration: other subjects in My-Subjects/ use semicolon-separated sources
% and include an "accessed" date for web documentation.

\newcommand{\OOPSUnitISources}{Schildt, \textit{Java: A Beginner's Guide} (10e), McGraw-Hill (Java fundamentals + OOP basics).; Oracle Java Tutorials: Getting Started + Language Basics + Classes and Objects (accessed \today).; Java SE API docs (e.g., \texttt{String}, \texttt{Scanner}) (accessed \today).}

\newcommand{\OOPSUnitIISources}{Schildt, \textit{Java: The Complete Reference} (12e), McGraw-Hill (classes, inheritance, packages, interfaces).; Bloch, \textit{Effective Java} (3e), Addison-Wesley, 2018 (design choices; interfaces vs abstract classes).; Oracle Java Tutorials: Classes and Objects + Interfaces + Packages (accessed \today).; Java SE API docs (e.g., \texttt{Comparable}, \texttt{Runnable}) (accessed \today).}

\newcommand{\OOPSUnitIIISources}{Schildt, \textit{Java: The Complete Reference} (12e) (nested/inner classes, exceptions, threads, I/O).; Oracle Java Tutorials: Nested Classes + Exceptions + Concurrency + Basic I/O (accessed \today).; Java SE API docs (e.g., \texttt{Thread}, \texttt{AutoCloseable}) (accessed \today).}

\newcommand{\OOPSUnitIVSources}{Schildt, \textit{Java: The Complete Reference} (12e) (generics, lambdas, Swing, JDBC).; Bloch, \textit{Effective Java} (3e) (generics + lambdas).; Oracle Java Tutorials: Generics + Lambda Expressions + Swing + JDBC Basics (accessed \today).; Java SE API docs (e.g., \texttt{java.util.function}, \texttt{javax.swing}, \texttt{java.sql}) (accessed \today).}

\newcommand{\OOPSUnitVSources}{Schildt, \textit{Java: The Complete Reference} (12e) (Collections Framework + wrapper classes).; Oracle Java docs: Collections Framework overview (accessed \today).; Java SE API docs (\texttt{java.util} collections; wrapper classes in \texttt{java.lang}) (accessed \today).}

\newcommand{\OOPSUnitVISources}{Git SCM: \textit{Pro Git} (2e) (version control workflow), accessed \today.; GitHub Docs (repositories, issues, pull requests), accessed \today.; Oracle Java Tutorials: Swing + JDBC (accessed \today).; JUnit 5 User Guide (accessed \today).; Apache Maven Project: Getting Started (accessed \today).; Gradle User Manual: Getting Started (accessed \today).; UML class diagrams: UML 2.x overview/spec (accessed \today).}

\newcommand{\OOPSMiscWhyInterfacesSources}{Schildt, \textit{Java: The Complete Reference} (12e) (interfaces).; Bloch, \textit{Effective Java} (3e) (prefer interfaces where appropriate).; Oracle Java Tutorials: Interfaces (accessed \today).; Java SE API docs (e.g., \texttt{Comparable}, \texttt{Runnable}) (accessed \today).}



\title[OOPS with Java]{Object Oriented Programming (CSEG1044)}
\subtitle{Unit 01 -- Lecture 01: OOP Paradigm + Principles + Modelling}
\begin{document}

\begin{frame}[fragile]
  \titlepage
  \vspace{0.5em}
  \slidesources{\OOPSUnitISources}
\end{frame}

\begin{frame}{Quick Links}
  \centering
  \hyperlink{sec:overview}{\beamerbutton{Overview}}\hspace{0.6em}
  \hyperlink{sec:concepts}{\beamerbutton{Key Topics}}\hspace{0.6em}
  \hyperlink{sec:demo}{\beamerbutton{Demo}}\hspace{0.6em}
  \hyperlink{sec:summary}{\beamerbutton{Summary}}
  \slidesources{\OOPSUnitISources}
\end{frame}

\begin{frame}{Agenda}
  \tableofcontents
  \slidesources{\OOPSUnitISources}
\end{frame}

\section{Overview}
\label{sec:overview}

\begin{frame}{Lecture Snapshot}
\begin{itemize}[<+->]
  \item \textbf{Unit focus:} Introduction to OOPs and Java Fundamentals
  \item \textbf{Course thread:} Use OOP Paradigm + Principles + Modelling to make the first text-based version of Campus Resource Manager easier to reason about before the full object model arrives.
\end{itemize}
  \slidesources{\OOPSUnitISources}
\end{frame}

\begin{frame}{Recall Before You Start}
\begin{itemize}[<+->]
  \item \textbf{Programming basics}: Variables, loops, and functions from C still matter; OOP changes how they are organized. Main reference: Programming in C prerequisite.
  \item \textbf{Problem decomposition}: OOP still depends on breaking problems into clear responsibilities before coding. Main reference: Programming in C prerequisite.
\end{itemize}
  \slidesources{\OOPSUnitISources}
\end{frame}


\begin{frame}{Learning Outcomes}
  \begin{itemize}[<+->]
    \item Explain why OOP became popular for large software systems.
    \item Define: abstraction, encapsulation, inheritance, polymorphism (with a small example).
    \item Convert a short problem statement into a small class design (classes + responsibilities).
  \end{itemize}
  \slidesources{\OOPSUnitISources}
\end{frame}

\begin{frame}{Key Terms to Know}
\small
\begin{columns}[t]
\column{0.48\textwidth}
\begin{itemize}
  \item \textbf{Class}: A blueprint (type) describing data (fields) and operations (methods).
  \item \textbf{Object}: A runtime instance of a class (has identity + state + behavior).
  \item \textbf{State}: The data inside an object (values of its fields at a time).
\end{itemize}
\column{0.48\textwidth}
\begin{itemize}
  \item \textbf{Behavior}: What an object can do (methods).
  \item \textbf{Invariant}: A rule that must always remain true for an object (e.g., \texttt{balance >= 0}).
  \item \textbf{Interface (idea)}: A public view/contract of what operations are allowed (we will formalize this later in Java).
\end{itemize}
\end{columns}
  \slidesources{\OOPSUnitISources}
\end{frame}

\begin{frame}{Study Flow for This Lecture}
\begin{enumerate}[<+->]
  \item Refresh the prerequisite bridge before reading new ideas in isolation.
  \item Study the main build in this order: \textit{Why OOP? (vs procedural)} $\rightarrow$ \textit{The 4 OOP principles (with intuition)} $\rightarrow$ \textit{How to model a real-world entity}.
  \item Trace the worked example and connect each line back to one concept frame.
  \item Attempt the mini exercises before moving to the recap and exit check.
\end{enumerate}
  \slidesources{\OOPSUnitISources}
\end{frame}


\section{Key Topics}
\label{sec:concepts}

\begin{frame}{Unit 01: Introduction to OOPs}
  \begin{itemize}[<+->]
    \item Lecture focus: OOP Paradigm + Principles + Modelling
  \end{itemize}
  \slidesources{\OOPSUnitISources}
\end{frame}

\begin{frame}{Today: Roadmap}
  \begin{itemize}[<+->]
    \item OOP history and evolution; why OOP (vs procedural)
    \item Principles: abstraction, encapsulation, inheritance, polymorphism
    \item Modelling a real-world entity (identify classes/objects/responsibilities)
  \end{itemize}
  \slidesources{\OOPSUnitISources}
\end{frame}

\begin{frame}{Why OOP? (motivation)}
  \begin{itemize}[<+->]
    \item Small programs: procedural style works well (few files, few developers).
    \item Large programs: managing shared state becomes difficult.
    \item OOP idea: \textbf{data + operations together} as an \textbf{object}.
    \item Benefit: better modularity, easier maintenance, clearer mapping to real-world entities.
  \end{itemize}
  \slidesources{\OOPSUnitISources}
\end{frame}

\begin{frame}{OOP building blocks}
  \begin{itemize}[<+->]
    \item \textbf{Class}: blueprint (fields + methods)
    \item \textbf{Object}: instance at runtime (identity + state + behavior)
    \item \textbf{State}: values of fields
    \item \textbf{Behavior}: methods that operate on the state
    \item \textbf{Invariant}: rule that must always hold (e.g., \texttt{balance >= 0})
  \end{itemize}
  \slidesources{\OOPSUnitISources}
\end{frame}

\begin{frame}{Principle 1: Abstraction}
  \begin{itemize}[<+->]
    \item \textbf{What}: focus on \emph{what} a thing does, not \emph{how} it does it.
    \item \textbf{Why}: hides complexity; makes code easier to use and change.
    \item \textbf{Example}: \texttt{BankAccount.deposit()} is visible; internal storage details are hidden.
    \item \textbf{Pitfall}: exposing internal details (public fields) breaks the model.
  \end{itemize}
  \slidesources{\OOPSUnitISources}
\end{frame}

\begin{frame}{Principle 2: Encapsulation}
  \begin{itemize}[<+->]
    \item \textbf{What}: keep data private; control changes via methods.
    \item \textbf{Why}: protects invariants; prevents accidental misuse.
    \item \textbf{Example}: only \texttt{deposit/withdraw} can change \texttt{balance}.
    \item \textbf{Pitfall}: setters that allow invalid states (e.g., \texttt{setBalance(-10)}).
  \end{itemize}
  \slidesources{\OOPSUnitISources}
\end{frame}

\begin{frame}{Principle 3: Inheritance}
  \begin{itemize}[<+->]
    \item \textbf{What}: models an \textbf{is-a} relationship (specialization).
    \item \textbf{Why}: reuse common behavior when relationship is truly is-a.
    \item \textbf{Example}: \texttt{SavingAccount} is-a \texttt{BankAccount}.
    \item \textbf{Pitfall}: using inheritance only for code reuse; prefer composition for has-a.
  \end{itemize}
  \slidesources{\OOPSUnitISources}
\end{frame}

\begin{frame}{Principle 4: Polymorphism}
  \begin{itemize}[<+->]
    \item \textbf{What}: one interface/type, many implementations.
    \item \textbf{Why}: enables flexible design (code works with parent types, not specific classes).
    \item \textbf{Example}: a \texttt{BankAccount} reference can point to different account subtypes.
    \item \textbf{Pitfall}: confusing \textbf{overloading} (compile-time) with \textbf{overriding} (runtime).
  \end{itemize}
  \slidesources{\OOPSUnitISources}
\end{frame}

\begin{frame}{Modelling: from problem statement to classes}
  \begin{enumerate}[<+->]
    \item List 3--5 \textbf{scenarios} (what users do).
    \item Extract \textbf{nouns} $\rightarrow$ candidate classes (Account, Customer).
    \item Extract \textbf{verbs} $\rightarrow$ candidate methods (deposit, withdraw).
    \item Assign \textbf{responsibilities} (avoid one ``God class'').
    \item Write \textbf{invariants} and validate with sample cases.
  \end{enumerate}
  \slidesources{\OOPSUnitISources}
\end{frame}

\begin{frame}{Checkpoint and Reflection}
\begin{block}{Reflection prompt}
Where would \textit{OOP Paradigm + Principles + Modelling} matter most in Campus Resource Manager, and which design choice would you justify using this lecture's ideas?
\end{block}
\begin{block}{Quick self-check}
Which part needs one more example before you move on: \textit{Why OOP? (vs procedural)}, \textit{The 4 OOP principles (with intuition)}, the worked example, or the recap?
\end{block}
  \slidesources{\OOPSUnitISources}
\end{frame}

\section{Demo / Example}
\label{sec:demo}

\begin{frame}[fragile,shrink=8]{Example: Model a real-world entity (1/2)}
\begin{lstlisting}
public class BankAccount {
    private final String accountNo;
    private double balance; // invariant: balance >= 0

    public BankAccount(String accountNo, double openingBalance) {
        if (openingBalance < 0) throw new IllegalArgumentException("openingBalance");
        this.accountNo = accountNo;
        this.balance = openingBalance;
    }

    public void deposit(double amount) {
        if (amount <= 0) throw new IllegalArgumentException("amount");
        balance += amount;
    }
\end{lstlisting}
  \slidesources{\OOPSUnitISources}
\end{frame}

\begin{frame}[fragile]{Example: Model a real-world entity (2/2)}
\begin{lstlisting}
// continued...
    public void withdraw(double amount) {
        if (amount <= 0 || amount > balance)
            throw new IllegalArgumentException("amount");
        balance -= amount;
    }
}
\end{lstlisting}
  \slidesources{\OOPSUnitISources}
\end{frame}

\begin{frame}{Mini Exercises (2--3 mins)}
  \begin{enumerate}[<+->]
    \item Model \texttt{Student}: 3 fields + 3 methods + 1 invariant.
    \item Decide: is-a or has-a?\; Car--Engine,\; Dog--Animal.
    \item Add \texttt{transferTo(other, amount)} to \texttt{BankAccount}.
  \end{enumerate}
  \slidesources{\OOPSUnitISources}
\end{frame}

\begin{frame}{Watch-outs}
\begin{itemize}[<+->]
  \item Treating a class as only a data bucket with no behavior.
  \item Confusing one object instance with the class blueprint.
  \item Using inheritance before the responsibilities are understood.
\end{itemize}
  \slidesources{\OOPSUnitISources}
\end{frame}

\section{Summary}
\label{sec:summary}

\begin{frame}{Key Takeaways}
  \begin{itemize}[<+->]
    \item OOP helps manage complexity by organizing code around objects (state + behavior).
    \item Abstraction + encapsulation simplify and protect designs.
    \item Inheritance + polymorphism enable reuse and flexibility (use carefully).
    \item Modelling is the bridge from real-world requirements to classes and responsibilities.
  \end{itemize}
  \slidesources{\OOPSUnitISources}
\end{frame}

\begin{frame}{Checkpoint Questions}
  \begin{enumerate}[<+->]
    \item Why do we keep fields private in OOP?
    \item Give one example each: abstraction and encapsulation.
    \item Give one example each: is-a and has-a.
  \end{enumerate}
  \slidesources{\OOPSUnitISources}
\end{frame}

\begin{frame}{Next Study Steps}
\begin{itemize}[<+->]
  \item Revisit the recall references if any older idea still feels weak.
  \item Rework the lecture example without looking at the notes first.
  \item Attempt the self-study practice and then answer the exit questions in full sentences.
  \item Apply the idea once inside Campus Resource Manager to make the concept stick.
\end{itemize}
  \slidesources{\OOPSUnitISources}
\end{frame}

\begin{frame}{Next Lecture Preview}
  \textbf{Next:} Java Overview + Program Structure + Features
  \vspace{0.6em}
  \begin{itemize}[<+->]
    \item Install/verify JDK; practice compile/run on your machine.
    \item Bring one question about JVM/JDK/JRE (confusing terms are expected!).
  \end{itemize}
  \slidesources{\OOPSUnitISources}
\end{frame}

\end{document}
